% Part: first-order-logic
% Chapter: tableaux
% Section: soundness.tex

\documentclass[../../../include/open-logic-section]{subfiles}

\begin{document}

\iftag{FOL}
      {\olfileid{fol}{tab}{sou}}
      {\olfileid{pl}{tab}{sou}}

\olsection{বিশুদ্ধতা}

\begin{explain}
ট্যাবলোর মতো কোনো !!^a{derivation} system তখনই \emph{বিশুদ্ধ}, যখন
তা এমন কিছু !!{derive} করতে পারে না যা প্রকৃতপক্ষে সত্য নয়। তাই
বিশুদ্ধতা !!{derivation} systems-এর এক ধরনের নিশ্চিত নিরাপত্তা-ধর্ম।
কোন প্রমাণতাত্ত্বিক ধর্মটি বিবেচ্য তার ওপর নির্ভর করে, যেমন আমরা চাই:
\begin{enumerate}
\item প্রতিটি !!{derivable}~$!A$ \iftag{FOL}{বৈধ}{সর্বতঃসত্য};
\item কোনো !!a{sentence} অন্য কয়েকটি থেকে !!{derivable} হলে সেটি
  তাদের যৌক্তিক ফলও;
\item কোনো !!{sentence}s-এর সমষ্টি অসঙ্গত হলে সেটি অপরিতৃপ্তিযোগ্য।
\end{enumerate}
এগুলি কোনো !!a{derivation} system-এর গুরুত্বপূর্ণ ধর্ম। এদের কোনোটি
না থাকলে !!{derivation} system-টিতে ত্রুটি আছে—তা অতিরিক্ত বিষয়
!!{derive} করবে। তাই কোনো !!a{derivation} system-এর বিশুদ্ধতা
প্রতিষ্ঠা করা সর্বাধিক গুরুত্বপূর্ণ।

এই সব প্রমাণতাত্ত্বিক ধারণা কোনো না কোনো ধরনের বদ্ধ !!{tableau}s
দিয়ে সংজ্ঞায়িত। তাই ওপরের (1)--(3) প্রমাণ করতে বদ্ধ !!{tableau}s-এর
অর্থতাত্ত্বিক ধর্ম সম্পর্কে কিছু প্রমাণ করা দরকার। প্রথমে কোনো
!!a{signed
  formula} একটি গঠনে পরিতৃপ্ত হওয়ার অর্থ সংজ্ঞায়িত করব।
তারপর দেখাব, কোনো !!{tableau} বদ্ধ হলে কোনো গঠনই তার সব অনুমিতি
পরিতৃপ্ত করে না। সেই ফল থেকে অনুসিদ্ধান্ত হিসেবে (1)--(3) পাওয়া যাবে।
\end{explain}

\begin{defn}
\iftag{FOL}{!!^a{structure}}{!!^a{valuation}}~
$\iftag{FOL}{\Struct{M}}{\pAssign{v}}$
!!a{signed formula}~$\sFmla{\True}{!A}$-কে \emph{পরিতৃপ্ত করে}
তখনই, যখন $\iftag{FOL}{\Sat{M}}{\pSat{v}}{!A}$; আর
$\sFmla{\False}{!A}$-কে পরিতৃপ্ত করে তখনই, যখন
$\iftag{FOL}{\Sat/{M}}{\pSat/{v}}{!A}$।
$\iftag{FOL}{\Struct{M}}{\pAssign{v}}$ কোনো
!!{signed formula}s-এর সমষ্টি~$\Gamma$-কে পরিতৃপ্ত করে তখনই, যখন
তা প্রতিটি $\sFmla{S}{!A} \in \Gamma$-কে পরিতৃপ্ত করে। সেটিকে
পরিতৃপ্ত করে এমন কোনো \iftag{FOL}{!!a{structure}}{!!a{valuation}}
থাকলে $\Gamma$~\emph{পরিতৃপ্তিযোগ্য}; অন্যথায়
\emph{অপরিতৃপ্তিযোগ্য}।
\end{defn}

\begin{thm}[বিশুদ্ধতা]
  \ollabel{thm:tableau-soundness}
  $\Gamma$-র একটি বদ্ধ !!{tableau} থাকলে $\Gamma$
  অপরিতৃপ্তিযোগ্য।
\end{thm}

\begin{proof}
কোনো !!a{tableau}-এর একটি শাখার !!{signed formula}s-এর সমষ্টি
পরিতৃপ্তিযোগ্য হলে এবং কেবল তখনই শাখাটিকে পরিতৃপ্তিযোগ্য বলি; আর
কোনো !!a{tableau}-তে অন্তত একটি পরিতৃপ্তিযোগ্য শাখা থাকলে সেটিকে
পরিতৃপ্তিযোগ্য বলি।

আমরা দেখাব: কোনো পরিতৃপ্তিযোগ্য !!{tableau}-কে অনুমান-বিধিগুলির
একটি দিয়ে প্রসারিত করলে সব সময় একটি পরিতৃপ্তিযোগ্য !!{tableau}
পাওয়া যায়। এতেই উপপাদ্যটি প্রমাণিত হবে। যেকোনো বদ্ধ !!{tableau}
শুধু~$\Gamma$-র অনুমিতি নিয়ে গঠিত !!{tableau}-তে অনুমান-বিধি
প্রয়োগ করে পাওয়া যায়। তাই $\Gamma$ পরিতৃপ্তিযোগ্য হলে তার জন্য
যেকোনো !!{tableau} পরিতৃপ্তিযোগ্য হতো। কিন্তু বদ্ধ !!{tableau}
স্পষ্টতই পরিতৃপ্তিযোগ্য নয়: প্রতিটি শাখায়
$\sFmla{\True}{!A}$ ও $\sFmla{\False}{!A}$ উভয়ই থাকে, আর কোনো গঠন
একই সঙ্গে~$!A$-কে পরিতৃপ্ত এবং অপরিতৃপ্ত করতে পারে না।

ধরা যাক আমাদের একটি পরিতৃপ্তিযোগ্য !!{tableau} আছে, অর্থাৎ এমন
!!a{tableau} যার অন্তত একটি পরিতৃপ্তিযোগ্য শাখা আছে। অনুমান-বিধি
প্রয়োগে হয় কোনো শাখায় !!{signed formula}s যোগ হয়, নয়তো কোনো
শাখা দুই ভাগে বিভক্ত হয়। বিধি-প্রয়োগে প্রসারিত হয়নি এমন কোনো
পরিতৃপ্তিযোগ্য শাখা যদি !!{tableau}-টিতে থাকে, তবে প্রসারিত
!!{tableau}-তেও সেটি পরিতৃপ্তিযোগ্য শাখা থেকে যায়; তাই প্রসারিত
ট্যাবলোটি পরিতৃপ্তিযোগ্য। সুতরাং শুধু সেই ক্ষেত্র বিবেচনা করলেই হবে,
যেখানে বিধিটি একটি পরিতৃপ্তিযোগ্য শাখায় প্রয়োগ করা হয়।

সেই শাখার !!{signed formula}s-এর সমষ্টি~$\Gamma$ এবং বিধিটি যে
!!{signed formula}-এ প্রয়োগ করা হয় তা
$\sFmla{S}{!A} \in \Gamma$ ধরি। বিধিটি শাখা ভাগ না করলে দেখাতে হবে
যে প্রসারিত শাখা—অর্থাৎ বিধিটির সিদ্ধান্তগুলি-সহ $\Gamma$—এখনও
পরিতৃপ্তিযোগ্য। বিধিটি শাখা ভাগ করলে দেখাতে হবে যে উৎপন্ন দুটি
শাখার অন্তত একটি পরিতৃপ্তিযোগ্য।

প্রথমে যেসব সম্ভাব্য অনুমানে শাখা ভাগ হয় না সেগুলি বিবেচনা করি।
\begin{enumerate}
\item $\sFmla{\True}{\lnot !B} \in \Gamma$-তে
  $\TRule{\True}{\lnot}$ প্রয়োগ করে শাখাটি প্রসারিত করা হলো। তখন
  প্রসারিত শাখায় !!{signed formula}s-এর সমষ্টি
  $\Gamma \cup \{\sFmla{\False}{!B}\}$ থাকে। ধরি
  $\iftag{FOL}{\Sat{M}}{\pSat{v}}{\Gamma}$। বিশেষত,
  $\iftag{FOL}{\Sat{M}}{\pSat{v}}{\lnot !B}$। অতএব
  $\iftag{FOL}{\Sat/{M}}{\pSat/{v}}{!B}$; অর্থাৎ
  $\iftag{FOL}{\Struct{M}}{\pAssign{v}}$
  $\sFmla{\False}{!B}$-কে পরিতৃপ্ত করে।
\item $\sFmla{\False}{\lnot !B} \in \Gamma$-তে
  $\TRule{\False}{\lnot}$ প্রয়োগ করে শাখা প্রসারিত করা হলো: অনুশীলন।
\item $\sFmla{\True}{!B \land !C} \in \Gamma$-তে
  $\TRule{\True}{\land}$ প্রয়োগ করে শাখা প্রসারিত করা হলো। এতে
  শাখায় দুটি নতুন !!{signed formula}s পাওয়া যায়:
  $\sFmla{\True}{!B}$ ও $\sFmla{\True}{!C}$। ধরি
  $\iftag{FOL}{\Sat{M}}{\pSat{v}}{\Gamma}$, বিশেষত
  $\iftag{FOL}{\Sat{M}}{\pSat{v}}{!B \land !C}$। তখন
  $\iftag{FOL}{\Sat{M}}{\pSat{v}}{!B}$ এবং
  $\iftag{FOL}{\Sat{M}}{\pSat{v}}{!C}$। অর্থাৎ
  $\iftag{FOL}{\Struct{M}}{\pAssign{v}}$
  $\sFmla{\True}{!B}$ ও $\sFmla{\True}{!C}$ উভয়কেই পরিতৃপ্ত করে।
\item $\sFmla{\False}{!B \lor !C} \in \Gamma$-তে
  $\TRule{\False}{\lor}$ প্রয়োগ করে শাখা প্রসারিত করা হলো: অনুশীলন।
\item $\sFmla{\False}{!B \lif !C} \in \Gamma$-তে
  $\TRule{\False}{\lif}$ প্রয়োগ করে শাখা প্রসারিত করা হলো। এতে
  শাখায় দুটি নতুন !!{signed formula}s পাওয়া যায়:
  $\sFmla{\True}{!B}$ ও $\sFmla{\False}{!C}$। ধরি
  $\iftag{FOL}{\Sat{M}}{\pSat{v}}{\Gamma}$, বিশেষত
  $\iftag{FOL}{\Sat/{M}}{\pSat/{v}}{!B \lif !C}$। তখন
  $\iftag{FOL}{\Sat{M}}{\pSat{v}}{!B}$ এবং
  $\iftag{FOL}{\Sat/{M}}{\pSat/{v}}{!C}$। অর্থাৎ
  $\iftag{FOL}{\Struct{M}}{\pAssign{v}}$
  $\sFmla{\True}{!B}$ ও $\sFmla{\False}{!C}$ উভয়কেই পরিতৃপ্ত করে।

\iftag{FOL}{%
\item $\sFmla{\True}{\lforall[x][!A(x)]} \in \Gamma$-তে
  $\TRule{\True}{\lforall}$ প্রয়োগ করে শাখা প্রসারিত করা হলো। এতে
  শাখায় নতুন !!{signed formula}~$\sFmla{\True}{!A(t)}$ পাওয়া যায়।
  ধরি $\Sat{M}{\Gamma}$, বিশেষত
  $\Sat{M}{\lforall[x][!A(x)]}$। \olref[syn][sem]{prop:quant-terms}
  অনুসারে $\Sat{M}{!A(t)}$। অতএব $\Struct{M}$
  $\sFmla{\True}{!A(t)}$-কে পরিতৃপ্ত করে।

\item $\sFmla{\False}{\lforall[x][!A(x)]} \in \Gamma$-তে
  $\TRule{\False}{\lforall}$ প্রয়োগ করে শাখা প্রসারিত করা হলো। এতে
  নতুন !!{signed formula}~$\sFmla{\False}{!A(a)}$ পাওয়া যায়, যেখানে
  $a$ এমন একটি !!a{constant} যা~$\Gamma$-তে নেই। $\Gamma$
  পরিতৃপ্তিযোগ্য বলে এমন একটি $\Struct{M}$ আছে যেন
  $\Sat{M}{\Gamma}$, বিশেষত
  $\Sat/{M}{\lforall[x][!A(x)]}$। দেখাতে হবে
  $\Gamma \cup \{\sFmla{\False}{!A(a)}\}$ পরিতৃপ্তিযোগ্য। এ জন্য
  নিচের মতো একটি উপযুক্ত~$\Struct{M'}$ সংজ্ঞায়িত করি।

  \olref[syn][ass]{prop:sat-quant} অনুসারে
  $\Sat/{M}{\lforall[x][!A(x)]}$ হবে তখনই, যখন কোনো $s$-এর জন্য
  $\Sat/{M}{!A(x)}[s]$। এবার $\Struct{M'}$-কে $\Struct{M}$-এর মতোই
  নিই, শুধু $\Assign{a}{M'} = s(x)$। যেহেতু $a$~$\Gamma$-তে নেই,
  তাই \olref[syn][ext]{cor:extensionality-sent} অনুসারে প্রতিটি
  $\sFmla{\True}{!C} \in \Gamma$-র জন্য $\Sat{M'}{!C}$ এবং প্রতিটি
  $\sFmla{\False}{!C} \in \Gamma$-র জন্য $\Sat/{M'}{!C}$।

  \olref[syn][ext]{prop:extensionality} অনুসারে
  $\Sat/{M'}{!A(x)}[s]$। \olref[syn][ext]{prop:ext-formulas} অনুসারে
  $\Sat/{M'}{!A(a)}[s]$। যেহেতু $!A(a)$ একটি !!a{sentence}, তাই
  \olref[syn][ass]{prop:sentence-sat-true} অনুসারে
  $\Sat/{M'}{!A(a)}$; অর্থাৎ $\Struct{M'}$
  $\sFmla{\False}{!A(a)}$-কে পরিতৃপ্ত করে।
\item $\sFmla{\True}{\lexists[x][!B(x)]} \in \Gamma$-তে
  $\TRule{\True}{\lexists}$ প্রয়োগ করে শাখা প্রসারিত করা হলো:
  অনুশীলন।
\item $\sFmla{\False}{\lexists[x][!B(x)]} \in \Gamma$-তে
  $\TRule{\False}{\lexists}$ প্রয়োগ করে শাখা প্রসারিত করা হলো:
  অনুশীলন।}{}
% BN-SRC-055 TARGET-CORRECTION-BEGIN
\par\smallskip\noindent\textnormal{\small\textbf{উৎস-সংশোধন (BN-SRC-055):} হিমায়িত ইংরেজি বিশুদ্ধতা-প্রমাণে সার্বিক পরিমাণসূচকের দুই ক্ষেত্রে একই বিধি-দৃষ্টান্তে $!A$ ও $!B$ অসঙ্গতভাবে মেশানো হয়েছে। সিদ্ধান্ত ও প্রতিস্থাপিত সূত্রের সঙ্গে মিলিয়ে বাংলা লক্ষ্যে সংশ্লিষ্ট পূর্বধারণা এবং অর্থতাত্ত্বিক ধাপগুলিতে $!A$ রাখা হয়েছে।} % BN-SRC-055 TARGET-CORRECTION-END
\end{enumerate}
এবার যেসব সম্ভাব্য অনুমানে শাখা ভাগ হয় সেগুলি বিবেচনা করি।
\begin{enumerate}
\item $\sFmla{\False}{!B \land !C} \in \Gamma$-তে
  $\TRule{\False}{\land}$ প্রয়োগ করে শাখা প্রসারিত করা হলো। এতে দুটি
  শাখা হয়: বাঁ শাখা $\sFmla{\False}{!B}$ দিয়ে এবং ডান শাখা
  $\sFmla{\False}{!C}$ দিয়ে এগোয়। ধরি
  $\iftag{FOL}{\Sat{M}}{\pSat{v}}{\Gamma}$, বিশেষত
  $\iftag{FOL}{\Sat/{M}}{\pSat/{v}}{!B \land !C}$। তখন
  $\iftag{FOL}{\Sat/{M}}{\pSat/{v}}{!B}$ অথবা
  $\iftag{FOL}{\Sat/{M}}{\pSat/{v}}{!C}$। প্রথম ক্ষেত্রে
  $\iftag{FOL}{\Struct{M}}{\pAssign{v}}$
  $\sFmla{\False}{!B}$-কে পরিতৃপ্ত করে, অর্থাৎ
  $\iftag{FOL}{\Struct{M}}{\pAssign{v}}$ বাঁ শাখার সূত্রগুলি
  পরিতৃপ্ত করে। দ্বিতীয় ক্ষেত্রে
  $\iftag{FOL}{\Struct{M}}{\pAssign{v}}$
  $\sFmla{\False}{!C}$-কে, অর্থাৎ
  $\iftag{FOL}{\Struct{M}}{\pAssign{v}}$ ডান শাখার সূত্রগুলি পরিতৃপ্ত করে।
\item $\sFmla{\True}{!B \lor !C} \in \Gamma$-তে
  $\TRule{\True}{\lor}$ প্রয়োগ করে শাখা প্রসারিত করা হলো: অনুশীলন।
\item $\sFmla{\True}{!B \lif !C} \in \Gamma$-তে
  $\TRule{\True}{\lif}$ প্রয়োগ করে শাখা প্রসারিত করা হলো: অনুশীলন।
\item \Cut{} প্রয়োগ করে শাখা প্রসারিত করা হলো। এতে দুটি শাখা হয়:
  একটিতে $\sFmla{\True}{!B}$, অন্যটিতে $\sFmla{\False}{!B}$।
  যেহেতু $\iftag{FOL}{\Sat{M}}{\pSat{v}}{\Gamma}$ এবং হয়
  $\iftag{FOL}{\Sat{M}}{\pSat{v}}{!B}$, নয়তো
  $\iftag{FOL}{\Sat/{M}}{\pSat/{v}}{!B}$, তাই
  $\iftag{FOL}{\Struct{M}}{\pAssign{v}}$ বাঁ বা ডান শাখার একটিকে
  পরিতৃপ্ত করে।
\end{enumerate}
\end{proof}

\tagprob{FOL}
\begin{prob}
\olref[fol][tab][sou]{thm:tableau-soundness}-এর প্রমাণ সম্পূর্ণ করো।
\end{prob}
\tagendprob

\tagprob{notFOL}
\begin{prob}
\olref[pl][tab][sou]{thm:tableau-soundness}-এর প্রমাণ সম্পূর্ণ করো।
\end{prob}
\tagendprob

\begin{cor}
\ollabel{cor:weak-soundness}
$\Proves !A$ হলে $!A$ \iftag{FOL}{বৈধ}{সর্বতঃসত্য}।
\end{cor}

\begin{cor}
\ollabel{cor:entailment-soundness}
$\Gamma \Proves !A$ হলে $\Gamma \Entails !A$।
\end{cor}

\begin{proof}
  $\Gamma \Proves !A$ হলে কোনো
  $!B_1$, \dots, $!B_n\in\Gamma$-র জন্য
  $\{\sFmla{\False}{!A}, \sFmla{\True}{!B_1}, \dots,
  \sFmla{\True}{!B_n}\}$-এর একটি বদ্ধ !!{tableau} থাকে।
  \olref{thm:tableau-soundness} অনুসারে প্রতিটি
  \iftag{FOL}{!!{structure}}{!!{valuation}}~
  $\iftag{FOL}{\Struct{M}}{\pAssign{v}}$ হয় কোনো $!B_i$-কে মিথ্যা
  করে, নয়তো $!A$-কে সত্য করে। তাই
  $\iftag{FOL}{\Sat{M}}{\pSat{v}}{\Gamma}$ হলে
  $\iftag{FOL}{\Sat{M}}{\pSat{v}}{!A}$-ও হয়।
\end{proof}

\begin{cor}
\ollabel{cor:consistency-soundness}
$\Gamma$ পরিতৃপ্তিযোগ্য হলে সেটি সঙ্গত।
\end{cor}

\begin{proof}
বিপরীত-প্রতিজ্ঞা প্রমাণ করি। ধরি $\Gamma$ সঙ্গত নয়। তখন
$!B_1$, \dots, $!B_n\in\Gamma$ এবং
$\{\sFmla{\True}{!B_1}, \dots, \sFmla{\True}{!B_n}\}$-এর একটি বদ্ধ
!!{tableau} আছে। \olref{thm:tableau-soundness} অনুসারে এমন কোনো
$\iftag{FOL}{\Struct{M}}{\pAssign{v}}$ নেই যেন প্রতিটি
$i=1$, \dots,~$n$-এর জন্য
$\iftag{FOL}{\Sat{M}}{\pSat{v}}{!B_i}$। কিন্তু তাহলে $\Gamma$
পরিতৃপ্তিযোগ্য নয়।
\end{proof}

\end{document}
